\section{Discussion and conclusions}\label{sec:conclusions}

Three points of the theory matter for solver design. First, pair hulls glued
on finitely many moments of a shared variable accept distributions that no
joint distribution reproduces, and they give no positive bound exactly when the
local zero sets alternate often enough. On three-variable paths, dense moment
relaxations with the missing product close these gaps, so there the advantage
of joint blocks over their pairs is one of representation, which support cuts
obtain in the model's own coordinates. For larger blocks this is no longer
true: on a four-variable path, and with a relative gap of nearly $100\%$ on the
full-gap instances of Proposition~\ref{prop:full-gap} with five and seven
variables, an exact support cut of the block is strictly stronger than the
dense first-level relaxation with all products. Second, for stars, the merged
blocks of Section~\ref{sec:composition}, the support problem remains exactly
solvable in nearly linear time when rows couple the center with one leaf at a
time. Third, as in safe-cut practice, certification does not constrain how
directions are found: only the final binary64 row, its bound over the whole
block domain and the row that the solver stores need to be checked, and all of
this can be replayed.

The computations add two findings to the theory. First, certification matters
in practice. The sample minimum would have given wrong constants for 38\% of
the MINLPLib cuts and for almost all sampled path-family cuts, and local
minimization still left wrong constants that cut off feasible solutions,
including known optima. A numerical global solve was a much better substitute
on these small blocks, but on one model it reported optimality for a constant
$1.8\cdot10^{-4}$ above the true minimum, which a fixed safety shift of
$10^{-6}$ does not cover. Second, the stored-row check exposed an obstacle that
the theory does not see: presolve aggregates or fixes the variables to which a
block refers, so a separator stated in the original variables loses cuts unless
it prevents these reductions, which weakens SCIP's own relaxation, or
certifies the floating-point presolve steps on which it relies.

On MINLPLib models the cuts did not help SCIP: they improved its root bound on
few models and changed no full solve, while SCIP's disabled-by-default
separators achieved as much or more on all but one of them. On the constructed
families, where the theory predicts a gap, the picture is different. With a
non-binding coupling row, SCIP's root bound stayed below the level of glued
pair relaxations in all three settings we ran, and only SCIP-nolocks came close
to it; the certified cuts went beyond this level when the separator could add
enough cuts per block and tried the exact support of each whole row first, and
this configuration solved every instance, against at most half of them for
each SCIP setting and for Gurobi. With a binding row, both direction rules
solved every instance, and the cuts that mattered came from the LP direction
search; but SCIP with its implicit-discreteness presolve step switched off
solved every instance as well, and faster, because its own relaxation already
captured most of the gap. On the star family, star blocks made the root bound
optimal to within $10^{-6}$, whereas blocks of four variables added almost
nothing to SCIP's bound; the solved count was then limited by SCIP's search for
good solutions, where Gurobi did better. Joint blocks larger than the
separator's default of four variables thus matter. In the solver the simpler
star oracle certified them in milliseconds, and the near-linear sweep of
Theorem~\ref{thm:star} handles stars with $10^4$ leaves in about two seconds.

Several questions remain open. Exact support for trees deeper than stars with
edge-local rows would extend the star algorithm to the structures of network
models; Remark~\ref{rem:trees} shows that the rational partition argument
does not extend directly. Proposition~\ref{prop:gain}(ii) reduces the gap of
the glued relaxation in any direction to a concave maximization over
re-splittings, but which directions to merge, and when merging pays off, is
not answered by it. A direction search with the guarantees of
Appendix~\ref{app:separation} and the speed of a solver callback is not
available. Finally, an activation rule that predicts when joint cuts pay off
would be needed before any default use; our structural rule did not achieve
this.
